\documentclass[11pt]{article}
% Shared preamble for the round 1 response letters.
% RED-D-26-00078, submitted as "Housing First or Treatment First?",
% resubmitted as "Housing First or Treatment First? Ask the Kids, Not the Parents"
%
% House style, decided once here:
%   * responses are in RED, referee text is in the black italic block above them
%   * NO HYPHENS anywhere in the prose
%   * every response ends with a pointer to where in the new paper it lives

\usepackage[margin=1.05in]{geometry}
\usepackage{amsmath,amssymb}
\usepackage{xcolor}
\usepackage[official]{eurosym}
\usepackage{booktabs}
\usepackage{longtable}
\usepackage{array}
\usepackage[hidelinks]{hyperref}

% enumitem is not in this TeX installation, so first level list spacing is set
% directly instead
\makeatletter
\renewcommand{\@listI}{%
  \leftmargin=1.7em \labelwidth=1.2em
  \topsep=0.4em \parsep=0pt \itemsep=0.3em \partopsep=0pt}
\makeatother

\definecolor{respred}{RGB}{190,25,25}

\setlength{\parskip}{0.35em}
\setlength{\parindent}{0pt}

\newcolumntype{L}[1]{>{\raggedright\arraybackslash}p{#1}}
\newcommand{\pname}[1]{\texttt{#1}}

% ---- the two building blocks -------------------------------------------------
% \cmt{label}{short title} opens a comment, then the verbatim or paraphrased
% referee text goes in the referee environment, then the response.

\newcommand{\cmt}[2]{%
  \par\medskip
  \noindent\textbf{#1\quad #2}\par}

\newenvironment{referee}
  {\begin{list}{}{\leftmargin=1.6em \rightmargin=1.2em \topsep=0.4em}\item[]\small\itshape}
  {\end{list}}

\newenvironment{response}
  {\par\begingroup\color{respred}\medskip}
  {\par\endgroup\medskip}

% a flag for statements in the reports that are factually wrong, so they are
% easy to find again and easy to handle carefully
\newcommand{\correction}{\textbf{[Point of correction.]}\ }
\newcommand{\concede}{\textbf{[Conceded.]}\ }
\newcommand{\newresult}{\textbf{[New result.]}\ }
\newcommand{\open}{\textbf{[Still open.]}\ }

% numbers that are not in yet, so they are impossible to miss on a read through
\newcommand{\tbd}[1]{\textbf{[TO FILL: #1]}}

% Marks what has been done since the first draft of these letters. The letters
% were written when most of this was a plan; this flag marks the places where it
% is now built and running.
\newcommand{\thisround}{\textbf{[Built since the first draft of this letter.]}\ }

\newcommand{\lettertitle}[2]{%
  \begin{center}
  {\Large\bfseries Housing First or Treatment First?}\\[0.3em]
  {\normalsize Manuscript RED-D-26-00078}\\[0.9em]
  {\large\bfseries #1}\\[0.3em]
  {\normalsize #2}
  \end{center}
  \vspace{0.6em}}


\begin{document}

\lettertitle{Response to the Associate Editor}{Comments A1 to A20, responses in red}

Thank you for a letter that was unusually specific about what a successful
revision would have to contain. The three headings you set out, careful modelling
of the mechanisms, careful modelling of the policies, and careful quantification,
are how I organised the revision.


{\color{respred}
\textbf{One note before the comments.} The paper has a new title: it was submitted as
\emph{Housing First or Treatment First?} and is now \emph{Housing First or Treatment First? Ask the Kids, Not the Parents}. The reason is the substance of the revision. The answer
now turns on the children, through the intergenerational linkage that your A6 and
Referee~2's comment C3 point to, and no longer on a loss of consortium added to
welfare after the model is solved (A9).

\textbf{What changed, under your three headings.} On the mechanisms, Section~3 is a
new two period model, calibrated, that isolates the one mechanism the rest of the
paper depends on: drug use concentrates at the bottom of the income distribution
because the mortality risk of a substance use disorder (SUD) cannot be insured,
while its unemployment risk can. The quantitative model drops the ingredients you
questioned in A2 to A7: labour supply, homeownership, house prices and rents, the
unemployment insurance (UI) benefit state and the type dependent transfers. It
also drops the loss of consortium, and the fiscal cost the previous draft charged
for each homeless household and each household with a SUD (A15). It adds detail where you asked for it, in the
joint process of homelessness, drug use and income: serious mental illness (SMI),
recovery that is harder on the street, treatment as a choice, an endogenous drug
price, and children who share their parent's consumption and dwelling and whose
adult productivity depends on their parent's SUD. No preference varies with any
state. On the policies, Treatment First now requires a household with a SUD to be
in treatment, and its drug test catches only some users, costs the government
money, and excludes those it catches for a time. On the quantification, every
internally calibrated parameter is reported with its data and model moment, the
26 parameters that carry the story are separated from the 33 standard ones, and
Section~5.2 validates the model on moments it was not fitted to, including the
two pieces of evidence you and Referee~2 pointed to, Costantini (2025) and Milby
et al.\ (2005).

The answer has changed with the model. Moral hazard alone does not justify
conditionality: with the externalities on children removed, a newborn is worse
off under Treatment First, by 1.09 percent of lifetime consumption, even when the
savings from the drug test buy vouchers for other households. With them, in the
main experiment, where the drug price and the probability of drawing a voucher
both clear, Treatment First cuts drug use by 21 percent and homelessness by 11
percent, and raises a newborn's welfare by 2.57 percent of lifetime consumption.
Three quarters of that gain is an externality on the child that its parent does
not count, and in general equilibrium the parents gain too, by 0.69 percent of
their own. The paper now also solves the transition, starting from the benchmark
distribution: neither homelessness nor the share of children in a shelter rises in
any year, and children born at the reform gain 1.45 percent.
}

%======================================================================
\section*{Careful modelling of the key mechanisms}

\cmt{A1}{A stylized model to convey the mechanisms}
\begin{referee}
``In order to convey the different mechanisms of the paper, it would be helpful
to provide a section with a simplified version of the model as suggested by both
R1 and R3. The sheer number of moving parts are currently obfuscating the
economics of the paper. \dots Derive basic model predictions on the relationships
between the key model outcomes (income, unemployment, housing, drug use) and show
these align with a few stylized facts.''
\end{referee}
\begin{response}
Done. Section~3 is a stylized model with two periods of thirty years each, a
continuum of households whose productivity is fixed over life and lognormal
across them, one asset, and a choice of whether and how much to use. Using carries
the risk of a SUD, and a SUD brings two risks: a lower chance of surviving into the
second period and a higher chance of being out of work in it. The model is dynamic
rather than static, which is the point of Referee~1's comment 1(e): the cost of a
SUD arrives in the second period, so the decision to use is a decision about a
continuation value, and a static model cannot represent it. It is also calibrated,
to the same United States data as the quantitative model (Tables~1 and~2), and its
five internally calibrated parameters hit their targets exactly, so its results are
numbers rather than signs.

It is built to show one thing, on which the rest of the paper relies: why drug use
concentrates at the bottom of the income distribution. Fewer than one in ten adults
have a SUD, mild included, and the average user's income is 0.74 of the population
average in the 2024 National Survey on Drug Use and Health (NSDUH). In the model
drugs are a luxury good, so without consequences use would concentrate at the top:
the average user would earn 3.51 times the population average. Figure~1 then turns
on each risk on its own. The risk of dying alone takes the average user to 0.64
times the average at its calibrated strength. The risk of unemployment alone leaves
it at 3.51 even at twice its calibrated strength, because a household can save
against it, so the loss is roughly proportional to its income. Death cannot be
insured against, and what it destroys grows with income. With saving switched off,
unemployment reverses the gradient as well, which confirms that insurability is
what separates the two. Footnote~10 reports the same result in a version with
housing, a shelter, UI and vouchers.

This motivates two decisions in the quantitative model, which keeps the stylized
model's preferences and its two risks. No preference varies with income, or with
any other state (A8): the income gradient comes from the mortality risk. And the
mortality multiplier of a SUD is calibrated to the income of users relative to the
average, because it is the parameter that decides how far down the income
distribution use sits (Block~1 of Section~5.1).

On the stylized facts: the two above open Section~3. Figure~2 in Section~5.2 shows
the share of households that use drugs at each level of income relative to the
mean, in the model and in NSDUH 2024, and use is highest at the bottom and lowest
at the top in both: 18.3 percent against 13.7 below 0.22 of the mean, and 4.2
against 4.8 above 1.62 times it. The relationships with unemployment and housing
are in the quantitative model, since the stylized model leaves housing out to
isolate the one mechanism. Unemployment among households with a SUD is 14.5 percent
in the model against 16.2 in NSDUH 2024 (Table~4, Block~2), and the shares of the
homeless with a SUD and with SMI are calibration targets (A13 and A16).

This is in Section~3 (Tables~1 and~2, Figure~1), with the facts in Section~5.2
(Figure~2 and Table~4).
\end{response}

\cmt{A2}{Is modelling unemployment insurance expiration important?}
\begin{referee}
``Is it important to model the expiration of unemployment insurance?''
\end{referee}
\begin{response}
\concede Not as a state of its own, and it is no longer one. In the previous draft
a household that drew UI kept it for its whole spell, which needed a separate
state. UI now lasts one year: a share of those who lose their job receive it in
their first year out of work, and after that no unemployed household receives
anything. This is folded into employment status, which takes three values:
employed, insured in the first year out of work, and uninsured (equation~13,
Section~4.3). Nothing about it has to be identified from data on drug use: the
replacement rate and the cap come from state law as the Department of Labor
compiles it, and the share insured from the Bureau of Labor Statistics (Blocks~9
and~12 of Appendix~A).

I keep UI in this simple form because the paper's mechanism contrasts the
unemployment risk of a SUD, which households can insure, with its mortality risk,
which they cannot, and UI is one of the ways the first is insured, together with
saving and food assistance. Because it stops after a year, the extra time out of
work that a SUD causes is spent without it, which is how a SUD's effect on finding
a job raises the chance of losing housing (Block~2 of Section~5.1).

This is in Section~4.3 and Appendix~A.
\end{response}

\cmt{A3}{The role of labour supply}
\begin{referee}
``What is the role of labor supply in this framework?''
\end{referee}
\begin{response}
\concede None in the new paper, because it is no longer in the model. Earlier
versions solved the model with endogenous labour supply, and with rents,
homeownership and house prices. These margins had a negligible effect on the
results, so I removed them (footnote~19, Section~4.4). An employed household now
supplies its labour inelastically (Section~4.1).

The employment margin the paper needs is kept, and it is the extensive one: whether
a household works at all. A SUD and SMI each make a job harder to find, and
homelessness makes a job both easier to lose and harder to find (equation~21,
Section~4.8). These are the margins with data behind them. The SUD and SMI
multipliers are calibrated to the shares of the homeless with each condition
(Blocks~2 and~3 of Section~5.1), the penalty of homelessness is set from Desmond
and Gershenson (2016), and the unemployment rate of households with a SUD is a
validation moment, 14.5 percent in the model against 16.2 in NSDUH 2024 (Table~4,
Block~2).

This also answers Referee~3's minor comment 1, which asked why the labour supply
elasticity plays no clearer role when the budget is balanced with labour taxes. The
budget is no longer balanced with taxes. The tax parameters are held at their
benchmark values in every experiment, and the probability of drawing a voucher,
$\pi_V$, adjusts so that net government revenue stays at its benchmark level
(Section~4.11). No experiment raises a labour tax, so there is no deadweight loss
of that kind for a labour supply elasticity to govern.

This is in Sections~4.1, 4.4 and~4.11.
\end{response}

\cmt{A4}{Why homeownership rather than homeless against housed?}
\begin{referee}
``Why do we need the role of homeownership, and not simply the difference between
homelessness and houses?''
\end{referee}
\begin{response}
\concede It is not needed, and it is gone, together with house prices
(footnote~19). The housing choice is now a shelter and three rental sizes, and the
rent per unit of housing is held at its benchmark value in every experiment
(Sections~4.4 and~4.5). The shelter is literal homelessness, sleeping rough or in
an emergency shelter, and its service flow is calibrated to the homelessness rate
(Block~2 of Section~5.1).

I keep three rental sizes rather than a single housed state because how much of its
income a poor household must spend to stay housed is what Treatment First puts at
stake when it conditions the voucher. The smallest size is calibrated to the share
of renters paying more than half their income in rent, 19.2 percent in the 2022
Survey of Consumer Finances (SCF) against 25.8 in the model, and the two larger
sizes are set, not calibrated, at four and eight times the smallest (Block~8 of
Appendix~A). With these, 58.7 percent of housed households spend more than thirty
percent of their income on rent, against 43.2 percent of renters in the SCF, and a
voucher holder pays 39.0 percent of its income in rent, against 31.2 in the
Picture of Subsidized Households of the Department of Housing and Urban
Development (HUD), because a holder that rents above the fair market rent pays the
difference itself (Table~4, Block~3). This also bears on Referee~2's detail comment
on the housing grid, which had seven points of which three were calibrated: it now
has four, of which two are calibrated.

This is in Sections~4.4 and~4.5, Block~2 of Section~5.1 and Block~8 of Appendix~A.
\end{response}

\cmt{A5}{Why full general equilibrium?}
\begin{referee}
``Why does the policy analysis require a full general equilibrium model? Many
long-run partial equilibrium outcomes appear nearly identical to the general
equilibrium ones. Do account for the changes in the government, and ideally for
the adjustment in drug prices (See R1 Comment 1c) not necessarily in terms of
housing prices and rents.''
\end{referee}
\begin{response}
Done as you specify. House prices and rents no longer clear: there is no
homeownership, and the rent is fixed (Section~4.4). Two objects clear instead. The
government holds its net revenue at the benchmark level through the probability of
drawing a voucher, $\pi_V$, with taxes unchanged (Section~4.11). And the drug price
clears a supply curve whose price elasticity is 1.46, the estimate of Mace et al.\
(2020) for cannabis in Washington State (Section~4.12, Block~6 of Section~5.1). I
know of no direct estimate for an illicit drug. Cannabis is the closest market I
found, because it is still a Schedule~I substance at the federal level, so its
suppliers operate under federal prohibition (footnote~31).

Both matter, which is why the policy analysis needs them, and Table~5 adds them
one at a time so that the reader can see by how much. With the externalities on
children and both held at their Housing First levels, Treatment First cuts drug
use by 43 percent, a newborn gains 2.82 percent of lifetime consumption, and the
parents lose 1.19 percent of theirs. Letting $\pi_V$ clear spends the screen's
savings on vouchers for other households, and the parents then gain 2.42 percent.
Letting the drug price clear as well brings the fall in drug use down to 21 percent,
because the price falls by 5 percent and draws some households back into use. In
the main experiment, with both clearing, a newborn gains 2.57 percent and the
parents 0.69 percent. The partial equilibrium fall in drug use is twice the general
equilibrium one, and the partial equilibrium gets the sign of the parents' welfare
wrong.

This is in Sections~4.11 and~4.12, with the results by step in Section~6.2 and
Table~5.
\end{response}

\cmt{A6}{Intergenerational transmission discussed without intergenerational linkages}
\begin{referee}
``Why is there a discussion on intergenerational transmission in Section 4.1.3,
without intergenerational linkages?''
\end{referee}
\begin{response}
\newresult You were right that the previous draft discussed intergenerational
transmission with nothing in the model behind it: productivity persisted across
generations, but nothing a parent did reached its child. Referee~2's comment C3
asked for the opposite fix, a real linkage, on the ground that parents care about
their children's outcomes. I built the linkage, and it now decides the paper's
answer, which is why the title changed.

There are three links. First, the child lives with its parent for its first 25
years and shares the parent's consumption and dwelling, but not the drug. Every
unit a parent spends on drugs, and every unit of earnings it loses while a SUD
keeps it unemployed, is consumption and housing the child goes without, and a
homeless parent takes its child into the shelter (Section~4.2). Second, a child
whose parent had a SUD at any point while it lived at home starts adult life with
its productivity lowered by $\delta_d$, for life, and since a child's productivity
is drawn around its parent's, the penalty carries down the lineage (equation~22,
Section~4.8). Third, the parent values its child, but imperfectly, with the weight
$\omega_c$, both in the flow of childhood and in what its disorder costs the
child's adult life (Section~4.10). The share it leaves out, $1-\omega_c$, is an
externality. This gives three externalities on the child: the budget, earnings and
childhood homelessness externalities.

What the parent's disorder costs the child's adult life enters the parent's problem
as a difference, $D(y,d)$ in equation~(25): the child's value with a parent who
never had the disorder, less its value with the parent's actual history. It is zero
for a parent who never had a SUD and a cost otherwise. Written as a level, $\omega_c$
times the child's value, it would be worth more the likelier the parent is to die,
and so would pay most to households with a SUD for reasons that have nothing to do
with the child (footnote~23). Because the charge depends on whether the parent ever
had the disorder, the disorder state takes three values: never, recovered and
current (Section~4.1).

Each link is disciplined by evidence. The penalty, $\delta_d=0.24$, is the earnings
loss at age 32 that Nilsson (2017) measures for children exposed in the womb to a
Swedish alcohol experiment; to my knowledge, it is the only natural experiment
estimate of how a parental substance shock reaches a child's adult earnings. It is
checked indirectly: in the model, a child of a parent with a SUD has 1.46 times the
odds of acquiring one, at the low end of the 1.44 to 2.51 in Mellentin et al.\
(2016) and Yoon et al.\ (2013), and with the penalty switched off the odds ratio
is 1.02 in the economy without the externalities, so the transmission runs
through $\delta_d$ (Table~4, Block~6). The weight
$\omega_c$ is calibrated, to 0.085, from behaviour. In NSDUH 2024, an unmarried
adult living with one other adult and no child is 17.5 percent more likely to have
used illicit drugs in the past year than a lone parent living with one child,
reweighted to the same age, sex, marital status and family income. $\omega_c$ is
set so that removing the parent's concern for its child raises drug use in the
model by as much, 17.7 percent (Block~5 of Section~5.1). And 5.8 percent of
children in the model live with a parent who has a SUD, against 4.6 percent in
NSDUH 2024.

How much rests on the linkage is shown directly. The first five columns of
Table~5 remove the children: the parent gives them no weight, their productivity
does not depend on the parent, and childhood is not counted. A newborn then loses
under Treatment First in every configuration, by 1.17 percent of lifetime
consumption with the drug price and $\pi_V$ held and by 0.74 percent with both
clearing. With the children, it gains 2.57 percent in the main experiment, and
Table~6 shows that three quarters of that is the externality: 1.50 percent from
childhood and 0.46 percent from better productivity at entry. With
$\omega_c=0.69$, the weight Nishiyama (2000) estimates and the upper end of the
literature, a newborn still gains, by 3.18 percent, but only 23 percent of the gain
is then an externality (footnote~28). Referee~2 expected the channel to amplify the
gains from reducing drug use. In this model it decides their sign.

This is in Sections~4.2, 4.8 and~4.10, Block~5 of Section~5.1, Block~6 of
Section~5.2, and Sections~6.2 and~6.3 (Tables~5 and~6).
\end{response}

\cmt{A7}{Lump sum transfers and type dependent bequests}
\begin{referee}
``Incidentally, how are the lump-sum transfers determined as a function of types?
And why is it important that the bequests are type-dependent?''
\end{referee}
\begin{response}
\concede It was not important, and both are gone. In the previous draft every
living household received a transfer each period from the estates of the dead, and
the transfer depended on its productivity. Now only a household entering adult life
receives anything: a one time endowment, $T_0$, the same for every entrant, paid out
of the estates of the households who die, and the government collects the rest of
the estates (Section~4.3 and equation~26). $T_0$ has a job in the calibration: it
decides how likely a young adult is to start out in a shelter, and it is calibrated
to the share of the homeless aged 25 to 34, 26.0 percent in HUD's 2024 point in
time (PiT) count against 26.4 in the model.

This is in Section~4.3 and Block~2 of Section~5.1.
\end{response}

\cmt{A2 to A7}{Simplify the secondary dimensions, and model the joint process in more detail}
\begin{referee}
``Each of these model ingredients seem to be of lower importance relative to
modeling of the joint process of homelessness, addiction, and income. I would
prefer to see more detailed modeling of the key interactions of homelessness,
addiction, and income, while simplifying the framework in secondary dimensions.
While one can include additional features in a more detailed model in the appendix
for sensitivity analysis, the paper would benefit from a more focused analysis in
the body of the paper.''
\end{referee}
\begin{response}
This is what the revision did. The secondary dimensions of A2 to A7 are cut or
reduced to their simplest form: labour supply, homeownership, house prices and
rents, the separate UI state and the type dependent transfers. So is one dimension
of the drug choice: the chance that use leads to a SUD no longer rises with how
much a household uses (footnote~20). Rather than carry the cut margins in an
appendix, I report in footnote~19 that they had a negligible effect on the results.

The joint process of homelessness, drug use and income is modelled in more detail,
and its links run both ways. A SUD makes a job harder to find, and a household out
of work is likelier to lose its housing. Homelessness in turn makes a job easier to
lose and harder to find, raises mortality, and lowers the chance of recovering from
a SUD (Sections~4.5, 4.6 and~4.8). SMI is a state of the model: a SUD makes it more
likely to start, and it makes a job harder to find, raises mortality and makes
recovery from a SUD less likely (Section~4.7). Treatment is a choice that raises
recovery (Section~4.6). And a parent's SUD reaches its child's consumption, housing
and adult productivity (A6). Each hazard is anchored on, or checked against, a
clinical or epidemiological estimate (Section~2 and Blocks~2 to~4 of Section~5.1).

This is in Section~4 and Section~5.1.
\end{response}

\cmt{A8}{Preferences: keep the weights invariant and vary $b$ instead}
\begin{referee}
``A parsimonious way that came to my mind is to keep the Cobb-Douglas weights
invariant to addiction but let the marginal utility vary by addiction by changing
the parameter b decrease by addiction status.''
\end{referee}
\begin{response}
I went further than your suggestion: no preference parameter varies with anything.
The Cobb-Douglas weights and the shifter $b$ are the same for every household in
every state, with or without a SUD or SMI, housed or homeless (equations~11
and~12, Section~4.2). SMI does not enter the utility function at all
(Section~4.7). The one term in flow utility beyond the parent's and the child's
bundles and the value of being alive is the utility cost of a year in treatment,
which a household pays only if it chooses treatment.

The reason is the concern behind your comment, and behind Referee~2's comment C5
and Referee~3's comment M2(a). Utility rises with $b$, so if $b$ fell with the
disorder, a household that acquired one would also have its utility from a given
bundle lowered by assumption. A policy that prevented disorders would then raise
welfare through that assumption as well as through behaviour, and that mechanical
part is what the three of you were worried about. With preferences fixed, a SUD
costs a household only through what it does to its budget and to the hazards it
faces: a higher chance of dying, a job that is harder to find, a higher chance of
SMI, and, for its child, a lower adult productivity. Each of these is set from, or
calibrated to, data (Section~5.1), so the welfare effects in Section~6 come from behaviour and these
costs alone.

Two consequences follow. First, the income gradient in drug use cannot come from
preferences that differ with income, which is how the literature has produced it
(Greenwood et al., 2025); it comes from the uninsurable mortality risk of a SUD, as
the stylized model shows (A1). Second, the effect of housing on drug use cannot
come from a preference either. It runs through the chance of recovering, which is
lower in a shelter (A13).

On robustness to preferences: the stylized model's results hold with log utility
(footnote~8), and the parent's weight on its child, the preference parameter that
decides how much of the gain is an externality, is varied in footnote~28 (A6).

This is in Sections~4.2 and~4.7.
\end{response}

\cmt{A9}{Loss of consortium is ad hoc; present welfare and mortality separately}
\begin{referee}
``The consortium adjustment modeling also seems very ad-hoc while at the same
time mattering a lot of the welfare ranking as noted also by the Referees. A
possible way to address this would be to instead present the policy consequences
along two different dimensions, the welfare consequences of equation (30) and the
mortality consequences.''
\end{referee}
\begin{response}
\concede The loss of consortium is gone. Referee~1's version of the objection, that
the truly conservative value for a parameter with no source is zero and that zero
reversed the ranking, was right. Nothing is now added to welfare after the model is
solved. Welfare is the newborn's expected lifetime value in the model (equations~30
and~31, Section~6.1), and a death is valued only for what it costs within the
model, chiefly the flow value of being alive, $\aleph$, which is calibrated to a
value of a statistical life (VSL) of 80 times average labour income, 81 in the model (Block~1
of Section~5.1).

The case for Treatment First now rests on something the model contains. With the
children removed, a newborn loses under Treatment First; with them, it gains, and
Table~6 splits the gain into the part its parent counts and the externality (A6).
Referee~2's other suggestion, a disamenity that a large homeless population imposes
on everyone else, is not in the benchmark. I solved the model with one, at the value
in the contingent valuation of Loubi\`ere et al.\ (2020), about 0.2 percent of
average earnings a year, and no long run gain in Table~5 moves by more than 0.14
percentage points (footnote~4).

\open The paper does not report the mortality consequences of the policies as a
separate dimension. Your suggestion was a way to keep an ad hoc valuation of deaths
out of the welfare ranking, and with the adjustment gone the welfare numbers no
longer contain one. The model's mortality is checked instead against the evidence on
losing housing: losing a voucher raises the three year mortality of households with
SMI by about 3.6 percentage points, against the 4.6 by which Housing First lowers it
among homeless veterans with mental illness in Costantini (2025) (Table~4,
Block~7; see A14).

This is in Section~6.1 and Block~1 of Section~5.1, with the disamenity in
footnote~4.
\end{response}

%======================================================================
\section*{Modelling of the policies}

\cmt{A10}{Treatment First actively treats; the model has no treatment}
\begin{referee}
``Modeling `Treatment First' as a voucher program that depends on sobriety seems
to omit key characteristics. As R1 Comment 1 discusses, the program actively aims
to cure substance use disorder before housing individuals. The policy itself has a
direct role in impacting the transition out of addiction.''
\end{referee}
\begin{response}
\concede You and Referee~1 (comment 1(b)) were right: a policy named Treatment
First contained no treatment, and recovery was random.

A household with a SUD now chooses each year whether to be in treatment
(Section~4.6). Treatment doubles its chance of recovering, which is how much
medication for opioid use disorder raises remission (Paul et al., 2024), and
randomised trials of agonist medication against placebo find at least as large an
effect on abstinence (Fudala et al., 2003). It costs the household utility, and that
cost is calibrated to treatment take up: 24.5 percent of adults with a SUD received
treatment in the past year in NSDUH 2024, against 24.9 in the model. The government
pays for treatment under both policies, \$9,426 a year, 8.4 percent of average
earnings, the cost of a residential treatment episode in Roebuck et al.\ (2003), so
treatment competes with vouchers in the budget (Block~4 of Section~5.1). Recovery
also depends on housing and on SMI: it is less likely in a shelter, by a multiplier
calibrated to Milby et al.\ (2005), and less likely with SMI.

Treatment First now requires treatment: a household with a SUD keeps its voucher
only if it is in treatment and does not use, as the July 2025 executive order
directs and as the guidance implementing it requires of drug free housing
(Section~4.9). In the main experiment, the share of households with a SUD who are in
treatment rises from 24.6 to 29.0 percent (Table~5). The policy's direct effect on
recovery is checked against Milby et al.\ (2005), the one trial that compares
housing contingent on abstinence with housing without conditions: in the first year
of Treatment First, voucher holders with a SUD recover 1.40 times as fast as housed
households with a SUD and no voucher, against a ratio of 1.26 in the trial's weeks
of abstinence (Table~4, Block~8). The results do not depend on treatment: when a
household with a SUD cannot keep its voucher by entering treatment, Treatment First
cuts drug use by 29 percent and homelessness by 11 percent, and a newborn gains at
least as much (footnote~38).

This is in Sections~4.6 and~4.9, Block~4 of Section~5.1 and Block~8 of
Section~5.2.
\end{response}

\cmt{A11}{The policy is unlikely to be perfectly enforceable}
\begin{referee}
``Moreover, the policy is unlikely perfectly enforceable. R1 Comment 4 suggests
additional features of `Treatment First' that are of first-order importance for
the policy evaluation.''
\end{referee}
\begin{response}
Done. Under Treatment First, voucher holders and households that draw a voucher are
tested, and one that fails the condition is caught only with some probability, so a
user can keep its voucher (Section~4.9). The detection rate is 20 percent a year,
set so that the screen deters use as much as the closest evidence on a drug test
with a severe penalty, the US military's random testing: Mehay and Pacula (1999)
find that service members became 61 percent less likely to have used marijuana in
the past year, and in the model 62 percent of the users who face the screen stop
when it is introduced. Each test costs the government \$35, 0.031 percent of
average earnings, the upper end of the estimates (Legislative Budget Board, 2013),
so testing competes with vouchers in the budget. A household that is caught loses
its voucher and is excluded from the programme, returning to the lottery with
probability one third a year. That is an expected exclusion of three years, the
period for which the Housing Choice Voucher programme already bars a household
evicted for drug related criminal activity.

Referee~1's other features in comment 4 are handled as follows. False positives are
not in the model: a household that does not use, and is in treatment if it has a
SUD, is never flagged, because a positive screen is confirmed in a laboratory
before it has consequences. The cost of complying, for a household with a SUD, is
the utility cost of treatment, which is calibrated to take up (A10). Stigma and
administrative burdens beyond the cost of the test are not in the model.

The results hold at half the detection rate. At 10 percent, 41 percent of the users
subject to the screen stop on impact, inside the 11 to 45 percent range that Mehay
and Pacula find once those who left the military are counted. In the long run,
Treatment First then cuts drug use by 14 percent and homelessness by 6 percent, and
a newborn gains 1.65 percent, against 2.57 at the calibrated rate (footnotes~33
and~35). The previous draft's assumption, that every user is caught, appears in the
alternative transitions (Table~7): it raises homelessness by up to 14 percent in the
first years, so how much the screen catches decides what happens along the way.

This is in Section~4.9, in the paragraph on the screen at the start of Section~6,
and in Section~6.5.
\end{response}

%======================================================================
\section*{Quantification}

\cmt{A12}{Model fit never shown; untargeted moments missed}
\begin{referee}
``Both the model quantification and model validation are well below the bar for
publication. The paper never shows the model fit of the quantitative model and
the model does not match up to non-targeted moments, casting doubt on any
counterfactual policy experiments.''
\end{referee}
\begin{response}
\concede You were right on both counts, and both are addressed. Table~3 reports,
for each of the 26 parameters that carry the story, the moment, its source, its
value in the data and in the model, and the parameter's value; Table~A1 in
Appendix~A does the same for the 33 standard parameters. Of the 59 parameters, 21
are calibrated internally, and every one of them has its data and model moment side
by side, so the reader can judge the fit. The two largest misses, in percentage
points, are stated in the text: users earn 62 percent of the average in the model against 74 percent in the
data (Block~1 of Section~5.1), and 25.8 percent of renters pay more than half their
income in rent against 19.2 percent (Block~8 of Appendix~A).

Validation is Section~5.2 and Table~4, in eight blocks, each of which says why it
matters for the experiments. The model does well on what the experiments turn on.
Drug use by income has the shape of the data (Figure~2). The odds that a child of a
parent with a SUD acquires one, 1.46, are within the range of the evidence.
Unemployment among households with a SUD is 14.5 percent against 16.2, and
households with SMI are 2.64 times as likely as others to use drugs, against 2.44.
And the model reproduces two pieces of evidence it was not fitted to: the mortality
cost of losing housing for households with SMI (Costantini, 2025) and the effect of
conditioning housing on abstinence (Milby et al., 2005), which A13 and A14 discuss.

The misses that Referee~1 identified in comment 2(c) are smaller but not gone, and
Table~4 shows them. The model's homeless are poorer and less often employed than in
the data: 21.4 percent are employed, against 46, and their income is 6.5 percent of
the average, against 19.1 in Meyer et al.\ (2024). They are also younger: 37.9
percent are 35 to 44, against 27.2, and 11.5 percent are 55 and over, against 25.9.
Users are younger too, with a mean age of 34.6 against 39.6. And the participation
elasticity of drug demand is $-1.32$, larger in absolute value than the estimates of
$-0.3$ to $-0.9$; since Treatment First lowers the drug price, a larger price
response works against it (Block~5).

This is in Section~5.1 (Table~3), Section~5.2 (Table~4 and Figure~2) and
Appendix~A (Table~A1).
\end{response}

\cmt{A13}{Key elasticities the model should match}
\begin{referee}
``With respect to the policy evaluations, the key elasticities are (1) the
elasticity of drug use to housing access (see R3) and (2) the elasticity of
homelessness to drug addiction. The model should highlight these key elasticities
in the analysis and show that the model is able to match them.''
\end{referee}
\begin{response}
Both are now in the calibration, with the trial evidence that Referee~2 pointed to,
and each has its own entry in Section~5.1 saying what it does to the two policies.

\textbf{The elasticity of drug use to housing access.} In the model, housing
reaches drug use through recovery: a household with a SUD is less likely to recover
in a shelter, by the multiplier $\mu_H$ (equation~18). As $\mu_H$ falls, housing
buys more recovery and Housing First does more to cure a SUD. It is calibrated to
Milby et al.\ (2005), who randomised homeless adults with a cocaine disorder in
treatment into no housing, housing contingent on abstinence, and housing without
conditions: housing improved recovery by 47.3 percent, a ratio of 1.47 in the data
against 1.45 in the model (Block~4 of Section~5.1). The trial's second contrast,
contingent against unconditional housing, is the closest evidence there is on what
Treatment First does to those it houses. It is not targeted, and the model comes
close: 1.40 against 1.26 (Table~4, Block~8). The response of drug use to the
condition itself is set by the detection rate, against the military's random
testing (A11), and its response to the drug price is a validation moment (Table~4,
Block~5).

\textbf{The elasticity of homelessness to a SUD.} A SUD makes a job harder to find,
by the divisor $\kappa^{ue}_d$, and a household out of work longer is likelier to
lose its housing: this is the link from the disorder to the street, and it shapes
how much Treatment First deters. $\kappa^{ue}_d$ is calibrated to the share of the
homeless with a SUD, 21.7 percent in Gutwinski et al.\ (2021), 20.2 in the model
(Block~2 of Section~5.1), and the unemployment rate it implies for households with a
SUD is a validation moment, 14.5 percent against 16.2 (Table~4, Block~2). The same
holds for SMI, whose divisor is calibrated to the SMI share of the homeless, 25
percent against 24.7 in the model. Table~5 reports the SUD share of the homeless in
every experiment.

Neither runs through preferences (A8). Housing reaches drug use through the hazards
a household faces, which is also what the trial measures: Milby et al.\ report
abstinence by arm.

This is in Blocks~2 and~4 of Section~5.1 and Blocks~2, 5 and~8 of Section~5.2.
\end{response}

\cmt{A14}{Connect to Costantini and the other empirical work}
\begin{referee}
``For example, Constantini (2025) finds for a sample of 300,000 Veterans that
Housing First reduces three-year mortality by 4.6 percentage points while
Treatment First has no long-term health effects, in contrast with the model
simulations of this paper.''
\end{referee}
\begin{response}
\correction I read the comparison differently, and the paper now makes it
explicitly. Costantini estimates the effect of each of two programmes of the
Department of Veterans Affairs against no housing intervention. His Housing First is
the voucher programme the department runs with HUD, and his Treatment First is short term
transitional housing run by local contractors. The model reproduces his Housing
First finding. In Block~7 of Section~5.2, with the drug price and $\pi_V$ held, I
compare households with SMI whom Treatment First's revocation puts on the street
with the same households had they kept their voucher: losing it raises their three
year mortality by about 3.6 percentage points, against the 4.6 by which enrolling in
Housing First lowers it among his homeless veterans with mental illness (Table~4).
The mortality cost of losing housing in the model is of the same order as the one he
measures.

\open The paper reports no model counterpart to his Treatment First estimate. His
estimate compares a stay in transitional housing with no programme at all. The
paper's Treatment First is a different object: the same permanent voucher as
Housing First, kept only on condition of sobriety and, for a household with a SUD,
of treatment. The model has no time limited programme to set against his finding.

Where his evidence and the paper overlap, they agree. Like the trials, his
estimates speak to the households that enrol, within a given environment. The model
agrees at that level: with the drug price and $\pi_V$ held at their Housing First
levels, as in a trial, 74 percent of the households that have a SUD when the reform
lands are worse off under Treatment First, and 67 percent still are when both clear
(Section~7). What reverses the comparison lies outside what such designs observe:
the households the condition deters from starting to use, the vouchers its savings
buy for others, and the children (Section~7).

On the other papers you name: Abramson (2026) is now cited in Section~2, on how
eviction policy changes who becomes homeless. \open Cohen (2024), Brounstein and
Wieselthier (2026) and Humphries et al.\ (2026) are not cited in the current draft.

On your final question in this paragraph, whether a SUD has direct effects on
productivity, see A20.

This is in Section~2, Block~7 of Section~5.2 (Table~4) and Section~7.
\end{response}

\cmt{A15}{The fiscal cost of homelessness looks high}
\begin{referee}
``The fiscal cost of homelessness for the United States is also estimated in
recent work of Abramson and van Nieuwerburgh (forthcoming). It would be valuable
to provide a comparison of their figures to the CEA estimates or provide
robustness analysis with respect to these estimates of the fiscal costs. The
current estimate appears high (see R3).''
\end{referee}
\begin{response}
The fiscal cost of homelessness is no longer a parameter of the model, and neither
is the fiscal cost of a SUD that the previous draft took from the Council of
Economic Advisers. A household in a shelter costs the government nothing beyond the
transfers any household with its income receives: the budget (equation~26,
Section~4.11) pays UI, social security, food assistance, vouchers, drug tests and
treatment, and has no item for the shelter. The only public cost tied to a SUD is
treatment, \$9,426 a year for each household that takes it up (Block~4 of
Section~5.1). So neither the previous draft's figures nor Abramson and Van
Nieuwerburgh's enter the paper.

Leaving these costs out does not favour Treatment First. In the main experiment,
Treatment First lowers homelessness by 11 percent and the share of households with a
SUD by 19 percent in the long run, and along the transition homelessness does not
rise in any year (Sections~6.2 and~6.4). At those allocations, a positive cost per
homeless household or per household with a SUD would add to what Treatment First
saves the budget. This also settles Referee~2's comment C6 that the previous
draft's cost of an unhoused household was high and exceeded the rent on the
smallest unit: in the new model it is zero.

\open For the same reason, the paper does not compare Abramson and Van
Nieuwerburgh's figures with those of the Council of Economic Advisers, and does not
cite their paper: there is no fiscal cost of homelessness in the model for either
figure to discipline.

This is in Section~4.11.
\end{response}

\cmt{A16}{Where does 37 percent of homeless with SUD come from?}
\begin{referee}
``It is not clear where the statistic that 37 percent of homeless people are
experiencing SUD comes from. This seems very high and will also depend greatly on
the definition of homelessness \dots At a minimum, it should be clarified where
this number is coming from and what definition of homelessness the author is
targeting.''
\end{referee}
\begin{response}
\concede You were right to be suspicious, and the number is no longer used. The
previous draft took it as the midpoint of a range of substance abuse rates among
homeless adults in New York City and Philadelphia (Kuhn and Culhane, 1998), a
measure of substance abuse in general that varied with the population and the
definition. The paper is about drugs, so the target is now the prevalence of drug
use disorders among homeless people in the meta analysis of Gutwinski et al.\
(2021), 21.7 percent. It calibrates how much harder a SUD makes it to find a job,
and the model gives 20.2 percent (Table~3, Block~2).

On the definition: the homeless state is literal homelessness, sleeping rough or in
an emergency shelter (Section~4.5). A household that doubles up is housed in the
model. The homelessness rate is pinned not to a point but to a band, 0.40 to 1.00
percent of households: the floor is HUD's 2023 PiT count, which counts people in
shelters and on the street on a single night, and the ceiling is the share homeless
at some point over a year, from Burt et al.\ (2001). The model gives 0.62 percent.
The other moments of the homeless come from HUD's 2024 PiT count, for their ages,
and from Meyer et al.\ (2024), for their employment and income (Table~4, Block~1).

This is in Section~4.5, Block~2 of Section~5.1 and Block~1 of Section~5.2.
\end{response}

%======================================================================
\section*{Writing}

\cmt{A17}{The calibration section reads as a long list}
\begin{referee}
``Both the model parameterization and the calibration are rather thin. The
section is monotonically written as a longlist of targets. \dots estimates appear
to come from different decades, and it is often not clear which data waves are
used (such as the SCF), making the paper impossible to replicate \dots there
seems to be no distinction between parameters that are standard in the literature,
versus those that are new to the paper. From the references, for example, it is
not clear that the 10-year recovery rates for alcoholism (Vaillant, 1995) would
apply to substance use disorder for the homeless.''
\end{referee}
\begin{response}
Rewritten, along the three lines you set out.

\textbf{Insight rather than a list.} Section~5.1 groups the 26 parameters that carry
the paper's story into six blocks by the mechanism they work through: who uses
drugs, the costs of homelessness and vouchers, the SMI chain, recovery and
treatment, the externalities, and the drug market. Each block opens with what it
decides for the policies, and each parameter's entry runs in the same order: what
the parameter controls, how it bears on Treatment First, Housing First or drug use,
what happens as it rises, and only then its value, its moment and its source. The
validation in Section~5.2 is organised the same way, each block saying why it
matters for the experiments.

\textbf{Standard against new.} The 33 parameters common to most life cycle models,
for preferences, the rental ladder, earnings and employment, survival and the onset
of SMI, interest and taxes, social security and programme rules, and the rest of the
drug market, are in Appendix~A with
Table~A1, separate from the 26 in Table~3. Every row of both tables says whether the
parameter is set externally or calibrated internally, and for the internal ones
gives the data and the model value.

\textbf{Vintages and replication.} The surveys are named with their waves: NSDUH
2024, the 2022 SCF, HUD's PiT counts of 2023 and 2024, the 2024 Consumer
Expenditure Survey, and the Current Population Survey over 2014 to 2024, among
others. Estimates
taken from studies keep the study's date, which is why some are older. The
replication package's \pname{data/DATA\_APPENDIX.md} documents the dataset, year,
universe and computation behind every estimate, and the package carries all the
datasets, so nothing has to be downloaded (Section~5).

\textbf{Vaillant.} \concede You are right that a ten year remission rate for
alcoholism is not obviously the remission rate of a drug use disorder among the
homeless. The paper still uses it, but only for the baseline: the annual chance that
a housed household with a SUD recovers, 0.025, is the hazard implied by Vaillant's
ten year remission rate (Block~4 of Section~5.1). It is no longer applied to the
homeless. Recovery in a shelter is that baseline scaled by $\mu_H$, which is
calibrated to Milby et al.\ (2005), whose participants were homeless adults with a
cocaine disorder in treatment, the population the parameter is about (A13).

This is in Section~5 and Appendix~A.
\end{response}

\cmt{A18}{Fewer counterfactuals in the body}
\begin{referee}
``In terms of the policy analysis, ideally there would be fewer counterfactual
analysis in the body of the paper to convey the main message. \dots The first and
third experiment in this case seem to be most relevant singling out what is the
better policy under rationed quantities and non-rationed settings.''
\end{referee}
\begin{response}
Done, with one difference from your suggestion that I state plainly below. The
paper now asks one question, whether the existing rationed programme should require
sobriety, and has three experiments (Section~6). The first removes the externalities
on children and asks whether moral hazard alone justifies Treatment First; it does
not. The second introduces Treatment First in the full model in steps, letting
first $\pi_V$ and then the drug price clear, so that each general equilibrium effect
is shown on its own; the last step is the main experiment. Both are in Table~5. The
third is the transition from Housing First to Treatment First, with both clearing
every period (Figure~3). The alternative transitions, in which households with a SUD
cannot stop using, take a short section and one table (Section~6.5, Table~7), and
the robustness runs are reported in a sentence or a footnote each: half the
detection rate, a higher weight on the child, no earnings penalty, no response
through treatment, and a homelessness disamenity.

\open The difference is that no experiment removes the rationing. Every experiment
keeps the programme rationed, so the previous draft's third experiment, Treatment
First against Housing First with universal vouchers, is not in the paper. The reason
is the question: the executive order and HUD's funding notice that motivate the
paper change the conditions attached to the assistance that exists, not its scale
(Section~1). The size of the programme still moves in the experiments, because
$\pi_V$ spends what the screen saves: in the main experiment it rises from 6.75 to
7.80 percent.

This is in Section~6.
\end{response}

\cmt{A19}{Language}
\begin{referee}
``The language of `addicted' seems incorrect. The author should modify throughout
something to suffering from substance use disorder.''
\end{referee}
\begin{response}
Done. The paper speaks of households that use drugs and of households with a SUD
throughout the text, the tables and the figures. The word survives only in the
titles of cited works.

This runs through the whole paper.
\end{response}

\cmt{A20}{Direct productivity effects of substance use disorder}
\begin{referee}
``Incidentally, are there direct productivity effects of substance use addiction
on productivity?''
\end{referee}
\begin{response}
\open Not as a direct effect on the user's own productivity: a SUD does not lower
the wage of a household that works. It lowers the user's earnings through
employment instead. An unemployed household with a SUD finds work at the baseline
rate divided by $\kappa^{ue}_d$, so its spells out of work last longer
(Section~4.8). How much this costs is checked against the data: the unemployment
rate of households with a SUD is 14.5 percent in the model against 16.2 in NSDUH
2024 (Table~4, Block~2).

The productivity effect the model does have is on the next generation, and it is
central to the results: a child whose parent had a SUD while it lived at home
starts adult life with its productivity lowered by $\delta_d=0.24$, for life (A6).
This is the earnings externality, worth 0.46 percent of a newborn's lifetime
consumption in the main experiment, 18 percent of the gain (Table~6 and
footnote~30).

This is in Section~4.8 and Block~5 of Section~5.1.
\end{response}

\end{document}
